\documentclass[10pt]{article}
\usepackage[margin=20mm]{geometry}
\usepackage{amsmath,amssymb,fontspec,unicode-math,xurl,hyperref}
\IfFontExistsTF{DejaVu Serif}{\setmainfont{DejaVu Serif}}{\setmainfont{Latin Modern Roman}}
\setmathfont{latinmodern-math.otf}
\setlength{\parindent}{0pt}\setlength{\parskip}{6pt}\sloppy
\begin{document}
Abstract: The deterministic-equivalent comparison underlying the principal no-free-lunch theorem survives the checks below. A separate central contribution, the width-bottlenecked scaling law, has an invalid regime inference in its printed derivation: \(P\gg N\) alone does not force the effective ridge to scale as \(N^{-\alpha}\). A noiseless example satisfying the stated spectral and sample-width conditions instead has risk of order \(N^{-3/4}\) where the displayed law gives \(N^{-1}\). This is a Major issue for the scaling derivation as written. It does not refute the ridge-optimized monotonicity theorems or establish an error in the numerical experiments. Scientific adjudication remains pending.\par
Author: AI-assisted technical review\par
Date:\par
Title: Initial Scientific Review of
No Free Lunch from Random Feature Ensembles:
Scaling Laws and Near-Optimality Conditions\par
\# Identity, materials, and scope\par
This review concerns paper 0123, I25-43492, by Benjamin S. Ruben, William L. Tong, Hamza Tahir Chaudhry, and Cengiz Pehlevan, in the supplied 27-page ICML 2025 proceedings PDF. Page numbers below are one-based PDF pages and agree with the printed numbering. The original and source-manifest hashes were checked; neither file was changed. The pinned repository is `benruben87/Random\_Feature\_Ensembles`, commit . Its theory functions and selected experiment notebook cells were read statically. No author code was executed.\par
The main text, mathematical appendices A–E, and supplementary figure captions were inspected. The review independently checks the fixed-effective-ridge comparisons, the leading zero-ridge overparameterized expansion, and the width-scaling inference. It does not certify Gaussian universality, concentration results imported from other papers, empirical risk values, or every plotted point. The paper itself distinguishes its deterministic equivalent from finite-sample random risk in Section 3 and Appendix B; the findings respect that distinction.\par
\# Model and checks that hold\par
Write \(d_i=\mathrm{Df}_i(\kappa)\), \(A=-\kappa^2\mathrm{tf}'_1(\kappa)\), and \(B=\kappa\mathrm{tf}_1(\kappa)\). Equations (12)–(20) give \[\rho=\frac{N-d_1}{N-d_2},\quad
 \gamma_1=\frac{(1-\rho)d_1+\rho d_2}{P},\quad
 \gamma_2=\frac{d_2}{P},\qquad
 E_K=B_z+\frac{E_1-B_z}{K},\] where \(B_z=(A+\sigma_\epsilon^2)/(1-\gamma_2)\) and \(E_1=[\rho A+(1-\rho)B+\sigma_\epsilon^2]/(1-\gamma_1)\). For positive effective ridge and a learnable task, \(B>A\geq0\) and \(d_1>d_2\). At fixed \(\kappa,P\), increasing \(N\) increases \(\rho\), decreases both the signal numerator and \(\gamma_1\), and leaves \(B_z\) fixed. Increasing \(P\) at fixed \(\kappa,N\) decreases both risk denominators’ inflation factors. These observations supply the numerator step implicit in Appendix C.1, as well as the stated monotonicity.\par
For the fixed-total-size comparison, put \(\nu=1/K\) and \(N=\nu M\). The decomposition in (C.15) has two nonpositive terms: the derivative in (C.7) is negative, and \[\nu(1-\rho)=\frac{\nu(d_1-d_2)}{\nu M-d_2}\] decreases with \(\nu\), while \(1-\gamma_1\) increases. Thus the core comparison in Theorem 5.1 is supported within this deterministic-equivalent model. This is not a theorem about every realization of a finite feature ensemble or about feature-learning networks.\par
There is also a useful exact check on Section 6.1. For \(\lambda=0\), \(N>P\), and \(\sigma_\epsilon=0\), the fixed point has \(d_1=P\), independent of \(N\). Direct substitution gives \[\operatorname{Var}_z=\frac{P\kappa\,\mathrm{tf}_1(\kappa)}{N-P}.\] Its leading term is \(P\kappa\,\mathrm{tf}_1(\kappa)/N\), agreeing with (D.21) and the \(1/(KN)\) correction in (24). No issue was found in that specific leading-order assertion.\par
\# Problem\par
Problem 1 (Missing ridge regime in the core width-scaling derivation).  Assessment: Major, scoped to the stated derivation and its unqualified scaling law. Section 6.2, equations (26)–(27), gives a risk exponent under \(N\ll P\) and source/capacity power laws. Appendix D, equations (D.1)–(D.6), retains a general ridge \(\lambda\). The first paragraph on page 25 then infers that the left side of (D.6) is large from \(P\gg N\), and concludes \(\kappa\sim N^{-\alpha}\). The displayed assumptions do not imply that inference.\par
\#\# Assumptions and decisive calculation\par
Work entirely inside the paper’s deterministic-equivalent equations, with \[\eta_t=t^{-2},\qquad \bar w_t=1,\qquad
 \sigma_\epsilon=0,\qquad K=1,\qquad
 P=N^{3/2},\qquad \lambda=1.\] One may take \(N=m^2\) to make \(P=m^3\) integral. The spectrum has infinite rank, finite trace, and \(\alpha=2\). The target has finite variance \(\sum_t t^{-2}\) and satisfies \(\bar w_t^2\eta_t=t^{-(1+2\alpha r)}\) with \(r=1/4\). In particular \(P/N\to\infty\), as required in the printed argument.\par
The exact fixed-point equation (14) can be rewritten as \[\kappa\left(1-\frac{d_1}{P}\right)
        \left(1-\frac{d_1}{N}\right)=\frac1P.
 \tag{1}\] For this spectrum, integral comparison gives \[d_1(\kappa)\sim \frac{\pi}{2\sqrt{\kappa}},
 \qquad
 d_2(\kappa)\sim \frac{\pi}{4\sqrt{\kappa}}.
 \tag{2}\] The right side of (1) bounds \(\kappa\) from below by \(P^{-1}\). Since \(d_1\) decreases, this implies \(d_1/N=O(\sqrt P/N)=O(N^{-1/4})\to0\), and also \(d_1/P\to0\). Returning to (1) therefore proves \[\kappa\sim P^{-1}=N^{-3/2},\] rather than \(N^{-2}\). This is a self-consistent asymptotic solution, not a claim based on a finite numerical fit.\par
Here \(\mathrm{tf}_1=d_1\), so \[A=-\kappa^2d_1'(\kappa)\sim\frac{\pi}{4}\sqrt{\kappa},
 \qquad
 B=\kappa d_1(\kappa)\sim\frac{\pi}{2}\sqrt{\kappa}.\] Furthermore \(\rho\to1\) and \(\gamma_1\to0\). Equation (12) thus yields \[E_1\sim \frac{\pi}{4}N^{-3/4}.
 \tag{3}\] In contrast, equation (27) at \(\ell=1,\alpha=2,r=1/4\) gives \(s=1\) and equation (26) gives \(E_1\sim N^{-1}\). The discrepancy is in the leading power, not a prefactor.\par
An independent scalar script solves (1), using the closed-form sum \(d_1=[(\pi/\sqrt\kappa)\coth(\pi/\sqrt\kappa)-1]/2\). For \(N=10^4,10^6,10^8\), it obtains respectively \(\kappa N^{3/2}=1.1714,1.0510,1.0158\) and \(E_1N^{3/4}=0.9172,0.8252,0.7978\), approaching the analytic limits \(1\) and \(\pi/4\). The script imports no author module. These numbers corroborate the proof; they are not new benchmark results.\par
\#\# Alternative reading and counterevidence\par
The strongest favorable reading is that Section 6.2 intends the noiseless, sufficiently small, or optimally tuned ridge regime throughout. That interpretation is consistent with the broader paper’s emphasis on optimized risk. It is also supported by the pinned experiment notebook: cells 48–50 use a ridge sweep and set \(\sigma_\epsilon=0\) for the synthetic scaling experiment. Figure 4 explicitly says optimal ridge. Under that strengthened reading, the counterexample with fixed \(\lambda=1\) is outside the intended scope. This review does not claim that optimized risk has exponent \(3/4\) or that a reported fitted exponent is wrong.\par
Nevertheless, the printed inference from \(P\gg N\) to the large-left-side claim is invalid even at zero label noise, and the displayed derivation does not state the necessary ridge restriction. The restriction must enter the mathematical claim, not merely a numerical configuration.\par
\#\# Dependency, correction, and bounds on impact\par
The inference supplies \(\kappa\sim N^{-\alpha}\), which is then substituted into (D.7)–(D.12) to obtain the central joint-growth exponent and near-optimality discussion. This is why the defect is assessed as Major for that derivation, despite its repair being potentially short. It does not depend on lower benchmark scores or failure of the other main theorem.\par
A sufficient regime for the intended width scaling is noiseless prediction with \(P/N\to\infty\) and \(\lambda/(PN^{-\alpha})\to0\); then the width constraint dominates the ridge constraint. A precise optimized-risk claim can instead specify its optimization and prove that the selected ridge enters the appropriate regime. A fixed positive noise floor must also be removed or accounted for before asserting decay of the full risk in equation (10). These are analytic scope corrections, not a verified revision of the entire paper. Boundary source exponents may introduce logarithmic factors, so exponent statements should distinguish leading log-slopes from strict asymptotic equivalence.\par
\# Overall assessment and unresolved checks\par
The supported result is one scoped Major mathematical derivation issue. Theorem 5.1 and the checked leading overparameterized expansion remain supported. The Gaussian-universality and concentration assumptions are acknowledged limitations in the original paper, not newly established failures. Numerical reproduction, historical correspondence of the pinned source to each figure, and a comprehensive errata search remain unverified.\par
Item | State\par
Scientific recommendation | Major; human decision pending\par
Original/source identity | Hash checked\par
Independent calculation | Completed, exit code 0\par
Author-code execution | None\par
Repair verification | Analytic correction only; no empirical rerun\par
Author adoption | Not sought; unknown\par
Remaining work | Resolve the stated scientific evidence gaps\par
Final TeX | Saved\par
Report PDF and visual inspection | See the final provenance record\par
Focused acceptance | Recorded separately in provenance; adjudication pending\par
9 B. S. Ruben, W. L. Tong, H. T. Chaudhry, and C. Pehlevan, No Free Lunch from Random Feature Ensembles: Scaling Laws and Near-Optimality Conditions, ICML 2025, supplied proceedings PDF, especially pp. 3–8 and 22–27. The same authors, Random\_Feature\_Ensembles, pinned source tree, commit ; `EnsRFTheory.py`, lines 7–17, 55–123, 127–183 and 214–250; `Experiments.ipynb`, cells 45–50 (zero-based). Stable proceedings record: https://proceedings.mlr.press/v267/ruben25a.html.\par
\# Evidence and provenance\par
Evidence is under `verification/0123/evidence`: original text extraction, identity check, static source excerpts, notebook-cell excerpt and hash, independent script/results/run log, and a rendered source page 25 inspected visually. The latter is source inspection, not inspection of a compiled review PDF. The initial Windows review lacked a located TeX compiler. Final compilation and report-PDF inspection are recorded separately in provenance.\par
Scientific work and report authorship were performed in the initial review chat. An auxiliary agent assisted only with contract and compiler discovery. No human review is recorded. The frozen task token is `[internal identifier omitted]`; task `[internal task]`, run `B002`, attempt `001`. Full hashes, coverage limits, and pending decisions are in the adjacent JSON files.\par
The subsequent focused acceptance rechecked the reported Problems against the preserved source passages, assumptions, calculations, counterevidence, and core dependencies. Original and acceptance roles, final artifact hashes, compilation, and actual report-page inspection are recorded separately in the companion provenance. Scientific adjudication remains pending and human reviews remain empty.\par
\end{document}
